\documentclass[11pt,a4paper]{article}

\usepackage[utf8]{inputenc}
\usepackage[T1]{fontenc}
\usepackage{lmodern}
\usepackage[brazil]{babel}
\usepackage[a4paper,margin=2.4cm]{geometry}
\usepackage{xcolor}
\usepackage{titlesec}
\usepackage{enumitem}
\usepackage{tabularx}
\usepackage{booktabs}
\usepackage{array}
\usepackage{fancyhdr}
\usepackage[colorlinks=true,linkcolor=verde,urlcolor=verde]{hyperref}

% Paleta (verde do app)
\definecolor{verde}{HTML}{047857}
\definecolor{verdeclaro}{HTML}{ECFDF5}
\definecolor{cinza}{HTML}{475569}
\definecolor{cinzaclaro}{HTML}{F1F5F9}

% Títulos
\titleformat{\section}{\Large\bfseries\color{verde}}{\thesection.}{0.5em}{}
\titleformat{\subsection}{\large\bfseries\color{cinza}}{\thesubsection}{0.5em}{}
\titlespacing*{\section}{0pt}{1.4em}{0.6em}

% Cabeçalho/rodapé
\pagestyle{fancy}
\fancyhf{}
\renewcommand{\headrulewidth}{0.4pt}
\renewcommand{\footrulewidth}{0.2pt}
\lhead{\small\color{cinza}Club Brésil — Bolão da Copa}
\rhead{\small\color{cinza}Documentação de Arquitetura}
\cfoot{\small\color{cinza}\thepage}

\newcommand{\caixa}[2]{%
  \par\medskip\noindent
  \colorbox{verdeclaro}{\begin{minipage}{\dimexpr\linewidth-2\fboxsep}
  \textbf{\color{verde}#1}\\[0.3em]#2
  \end{minipage}}\par\medskip}

\newcolumntype{L}{>{\raggedright\arraybackslash}X}

\setlist{itemsep=2pt,topsep=3pt}

\begin{document}

% ---------------- CAPA ----------------
\begin{titlepage}
\centering
\vspace*{3cm}
{\Huge\bfseries\color{verde} Club Brésil}\\[0.3cm]
{\Large Bolão da Copa do Mundo 2026}\\[2cm]
{\huge\bfseries Documentação de Arquitetura}\\[0.5cm]
{\large Guia completo de tecnologia, hospedagem e funcionamento}\\[3cm]
\begin{minipage}{0.8\textwidth}
\centering\color{cinza}
\large Documento de referência e estudo para apresentar o projeto:
o que foi construído, com quais tecnologias, onde está hospedado e por quê.
\end{minipage}
\vfill
{\large Enzo Vendramin}\\[0.2cm]
{\color{cinza} Junho de 2026}
\end{titlepage}

\tableofcontents
\newpage

% ---------------- 1 ----------------
\section{Visão geral}

O \textbf{Club Brésil} é uma plataforma web privada de bolão para a Copa do
Mundo de 2026, pensada para grupos pequenos (cerca de 20 a 50 participantes).
Substitui planilhas, formulários e mensagens soltas por um produto único,
moderno e simples, no qual cada participante:

\begin{itemize}
  \item cria uma conta e aguarda aprovação do administrador;
  \item registra palpites de placar para os jogos;
  \item acompanha sua pontuação, o ranking geral e o pódio.
\end{itemize}

Há um único \textbf{administrador} (você), que aprova participantes, cadastra
jogos, lança resultados e gerencia o sistema. Existem hoje \textbf{dois bolões
independentes} rodando com o mesmo código: um para amigos na França e outro
para amigos no Brasil (detalhes na Seção~\ref{sec:doisboloes}).

\caixa{Princípio de projeto}{Toda decisão priorizou simplicidade, clareza e
velocidade de uso no celular, mesmo para pessoas sem familiaridade com
tecnologia. ``Menos é mais'' foi a régua de cada escolha.}

% ---------------- 2 ----------------
\section{Tecnologias e linguagens}

O projeto usa uma stack moderna e enxuta, na qual \textbf{um único programa}
cuida tanto da tela (\textit{frontend}) quanto da lógica de servidor
(\textit{backend}). Isso reduz custo, manutenção e complexidade.

\begin{center}
\begin{tabularx}{\textwidth}{@{}>{\bfseries}p{3.6cm} L@{}}
\toprule
Camada & Tecnologia e por quê \\
\midrule
Linguagem & \textbf{TypeScript} --- JavaScript com tipos. Dá segurança às
regras de pontuação e aos formulários, evitando erros bobos. \\[4pt]
Framework & \textbf{Next.js 15} (App Router) sobre \textbf{React 19}. Um só
projeto para frontend e backend, com renderização no servidor (rápido no
celular) e \textit{Server Actions} (formulários enviam direto para funções de
servidor, sem precisar de uma API separada). \\[4pt]
Estilo & \textbf{Tailwind CSS 4}. Design consistente e responsivo sem
bibliotecas pesadas de componentes. \\[4pt]
Banco de dados & \textbf{PostgreSQL} (gerenciado pela Neon) acessado via
\textbf{Prisma ORM}. O Prisma traduz o banco em código com tipos, e permitiu
começar em SQLite no desenvolvimento e migrar para PostgreSQL na produção
mudando praticamente uma linha. \\[4pt]
Autenticação & Sistema próprio com \textbf{JWT} em cookie (biblioteca
\texttt{jose}) e senhas protegidas por \textbf{bcrypt}. Sem serviços externos,
sem custo, sem a complexidade do ``login com Google''. \\[4pt]
Bandeiras & Imagens do serviço gratuito \textit{flagcdn.com}. Aparecem iguais
em qualquer aparelho (emojis de bandeira não funcionam no Windows). \\
\bottomrule
\end{tabularx}
\end{center}

\caixa{Por que ``tudo em um'' (Next.js)?}{Em vez de manter um site e um servidor
de API separados, o Next.js une os dois. Resultado: um único artefato para
publicar, menos peças que podem quebrar e um aplicativo muito leve
(aproximadamente 100~kB de JavaScript no navegador), o que o torna rápido até
em conexões ruins de celular.}

% ---------------- 3 ----------------
\section{Arquitetura geral}

O caminho de uma ação do usuário, do toque na tela até o banco de dados:

\begin{center}
\colorbox{cinzaclaro}{\begin{minipage}{\dimexpr\linewidth-2\fboxsep}
\centering\vspace{4pt}
\textbf{Celular/navegador} $\;\longrightarrow\;$ \textbf{Vercel} (executa o
Next.js na região mais próxima) $\;\longrightarrow\;$ \textbf{Server Action}
(valida permissões e regras) $\;\longrightarrow\;$ \textbf{Prisma}
$\;\longrightarrow\;$ \textbf{PostgreSQL/Neon}
\vspace{4pt}
\end{minipage}}
\end{center}

\noindent Pontos-chave da arquitetura:

\begin{itemize}
  \item \textbf{Validação sempre no servidor.} Nenhuma regra depende do
  navegador: aprovação de usuário, papel de administrador, bloqueio de
  palpites e pontuação são verificados no servidor a cada ação. O cliente
  nunca é ``confiável''.
  \item \textbf{Status do jogo é calculado, nunca armazenado.} ``Aberto'',
  ``Fechado'' e ``Finalizado'' derivam do horário atual e do placar. Assim, o
  bloqueio de palpites acontece \textbf{exatamente} no início da partida, sem
  precisar de nenhuma tarefa agendada.
  \item \textbf{Pontuação materializada.} Quando o administrador lança um
  resultado, o sistema calcula e grava os pontos de cada palpite e recalcula o
  ranking. As telas de leitura ficam, então, triviais e rápidas.
\end{itemize}

% ---------------- 4 ----------------
\section{Modelo de dados}

Cinco entidades organizam tudo, sempre separadas por edição (Copa 2026, e
futuras 2030, 2034\ldots):

\begin{center}
\begin{tabularx}{\textwidth}{@{}>{\bfseries}p{2.8cm} L@{}}
\toprule
Entidade & O que guarda \\
\midrule
User & Nome (exibido no ranking), nome de usuário (login), senha protegida,
papel (administrador ou participante) e situação (pendente ou aprovado). \\[4pt]
Edition & Uma edição do bolão (ex.: Copa 2026), a seleção campeã e a
configuração de trava do ``chute do campeão''. \\[4pt]
Match & Um jogo: as duas seleções, data/hora (em UTC), fase, placar oficial e
o controle de liberação de palpites. \\[4pt]
Prediction & O palpite de um usuário em um jogo (gols A e B) e os pontos já
calculados. Único por usuário e jogo. \\[4pt]
Participation & Liga um usuário a uma edição e guarda seus totais: pontos,
placares exatos, acertos, posição atual, posição anterior (para a evolução) e
o chute do campeão. \\
\bottomrule
\end{tabularx}
\end{center}

% ---------------- 5 ----------------
\section{Regras de negócio}

\subsection{Pontuação}
A pontuação é \textbf{configurável por bolão} (ver Seção~\ref{sec:doisboloes}).
A regra geral:

\begin{itemize}
  \item \textbf{Placar exato} (acertou os gols dos dois lados): pontuação cheia;
  \item \textbf{Acertou o resultado} (quem venceu, ou empate, sem o placar
  exato): pontuação menor;
  \item \textbf{Errou:} zero;
  \item \textbf{Chute do campeão} (acertar quem levanta a taça): bônus
  somado ao ranking quando o administrador define a campeã.
\end{itemize}

As duas primeiras regras \textbf{não somam}: quem acerta o placar exato recebe
a pontuação cheia, que já inclui ter acertado o vencedor.

\caixa{Mata-mata e pênaltis}{Conta apenas o placar do tempo regulamentar mais
prorrogação. Disputa por pênaltis não altera o placar para fins de pontuação:
um $1\times1$ decidido nos pênaltis continua sendo empate.}

\subsection{Bloqueio e edição de palpites}
Cada jogo abre para palpites automaticamente alguns dias antes (evitando uma
tela enorme de uma vez) e \textbf{trava no horário exato do apito inicial}.
Antes disso, o participante pode editar livremente, a qualquer momento e aos
poucos. O administrador também pode liberar ou bloquear manualmente cada jogo.

\subsection{Visibilidade dos palpites}
Antes do início da partida, ninguém vê o palpite dos outros. A partir do apito
inicial, os palpites de todos ficam visíveis (na aba Histórico, em um clique
expansível) --- e assim permanecem após o fim do jogo, aumentando a
transparência.

\subsection{Ranking e empates}
O ranking é recalculado automaticamente a cada resultado lançado. Em caso de
empate de pontos, ambos ocupam a \textbf{mesma posição} (ex.: 1º, 1º, 3º), sem
critérios de desempate. O sistema guarda a posição anterior para mostrar a
evolução ($\uparrow$, $\downarrow$ ou estável) a cada rodada.

% ---------------- 6 ----------------
\section{Estrutura de telas}

\begin{center}
\begin{tabularx}{\textwidth}{@{}>{\ttfamily}p{2.6cm} L@{}}
\toprule
\normalfont\bfseries Rota & Função \\
\midrule
/ & Página pública (sem login): pódio, ranking e resultados. \\
/entrar, /cadastro & Login e criação de conta. \\
/aguardando & Tela de espera até a aprovação do administrador. \\
/inicio & Próximo jogo, palpites pendentes, pódio e ranking resumido. \\
/palpites & Jogos abertos, agrupados por dia; ``Salvar todos os palpites''. \\
/ranking & Pódio e tabela completa com evolução. \\
/historico & Jogos iniciados, seu palpite, pontos e palpites de todos. \\
/agenda & Todos os jogos da Copa, dia a dia (inclusive os que não valem
palpite), como consulta. \\
/senha & Trocar a própria senha. \\
/admin/* & Painel: participantes, jogos, resultados e visão de palpites. \\
\bottomrule
\end{tabularx}
\end{center}

\noindent A navegação principal é uma barra inferior fixa (padrão de aplicativo
nativo), com campos e botões grandes, teclado numérico automático nos placares
e respeito à área segura do iPhone. Tudo \textbf{mobile-first}.

% ---------------- 7 ----------------
\section{Autenticação e segurança}

\begin{itemize}
  \item \textbf{Login} por nome de usuário e senha; senhas guardadas como
  \textit{hash} bcrypt (nunca em texto puro).
  \item \textbf{Sessão} em cookie \texttt{httpOnly} assinado (JWT), válido por
  30 dias --- não acessível por scripts no navegador.
  \item \textbf{Aprovação manual:} novos cadastros ficam pendentes até o
  administrador aprovar.
  \item \textbf{Limite de tentativas} de login (proteção contra adivinhação de
  senha) e \textbf{cabeçalhos de segurança} (anti-\textit{clickjacking}, etc.).
  \item \textbf{HTTPS} automático em produção.
  \item A página pública não expõe dados privados --- apenas nome e pontuação.
\end{itemize}

% ---------------- 8 ----------------
\section{Hospedagem e infraestrutura}\label{sec:hosting}

Três serviços, todos no \textbf{plano gratuito}:

\begin{center}
\begin{tabularx}{\textwidth}{@{}>{\bfseries}p{2.6cm} p{3.2cm} L@{}}
\toprule
Serviço & Papel & Por que esse \\
\midrule
GitHub & Guarda o código-fonte (repositório privado). & Padrão da indústria;
integra-se à Vercel para publicar automaticamente. \\[4pt]
Vercel & Hospeda e executa o aplicativo. & É a criadora do Next.js (encaixe
perfeito), com HTTPS e rede global automáticos. Publica sozinha a cada
atualização do código. \\[4pt]
Neon & Banco de dados PostgreSQL gerenciado. & PostgreSQL ``sem servidor'',
com plano gratuito generoso e backups; funciona bem com a Vercel. \\
\bottomrule
\end{tabularx}
\end{center}

\caixa{Por que não foi usado SQLite em produção?}{No desenvolvimento, o banco
era um simples arquivo (SQLite) --- ótimo para testar localmente. Mas a Vercel
executa o código em ambiente ``sem servidor'', onde não há um disco fixo para
guardar esse arquivo. Por isso a produção usa PostgreSQL (Neon), um banco de
verdade acessível pela rede. A troca custou praticamente uma linha de
configuração, graças ao Prisma.}

\caixa{A lição da latência}{No primeiro deploy, o código rodava nos EUA e o
banco na Europa: cada tela cruzava o Atlântico várias vezes e demorava
$\approx$2,2~s. Ao fixar a execução na \textbf{mesma região} do banco e dos
usuários, o tempo caiu para $\approx$0,15~s --- cerca de 10$\times$ mais rápido.
Daí a importância de código e banco ficarem geograficamente juntos.}

% ---------------- 9 ----------------
\section{Os dois bolões: França e Brasil}\label{sec:doisboloes}

O \textbf{mesmo código} alimenta dois bolões totalmente independentes (dados,
participantes e rankings separados). As diferenças vêm apenas de configuração,
não de código duplicado:

\begin{center}
\begin{tabularx}{\textwidth}{@{}p{3.2cm} L L@{}}
\toprule
& \textbf{Bolão França} & \textbf{Bolão Brasil} \\
\midrule
Servidor (Vercel) & Frankfurt (\texttt{fra1}) & São Paulo (\texttt{gru1}) \\
Banco (Neon) & Frankfurt (UE) & São Paulo (BR) \\
Fuso exibido & Horário de Paris & Horário de Brasília \\
Placar exato & 3 pontos & 5 pontos \\
Acertar resultado & 1 ponto & 2 pontos \\
Chute do campeão & 10 pontos & 10 pontos \\
Web Analytics & Ativado & Desativado \\
\bottomrule
\end{tabularx}
\end{center}

\noindent Tecnicamente, cada bolão é um \textit{projeto} separado na Vercel,
ligado a uma \textit{branch} do mesmo repositório (\texttt{main} para a França,
\texttt{brasil} para o Brasil) e a um banco Neon próprio. Comportamentos como
fuso horário e pontuação são lidos de \textbf{variáveis de ambiente}
(\texttt{APP\_TZ}, \texttt{SCORE\_EXACT}, \texttt{SCORE\_OUTCOME},
\texttt{NEXT\_PUBLIC\_ANALYTICS}) --- por isso o mesmo código se comporta
diferente em cada lugar, sem precisar mantê-los à parte.

\caixa{Um não atrapalha o outro}{Por serem projetos e bancos separados, os dois
bolões não compartilham servidor nem dados. Carga, lentidão ou qualquer
problema em um \textbf{não afeta} o outro.}

% ---------------- 10 ----------------
\section{Fluxo de atualização (publicação)}

O desenvolvimento e a publicação seguem um ciclo simples e seguro:

\begin{enumerate}
  \item A alteração é feita e testada \textbf{localmente} (no computador), sem
  afetar os sites no ar.
  \item Ao enviar o código para o GitHub (\texttt{git push}), a Vercel
  \textbf{detecta e publica} a nova versão em $\approx$2~minutos,
  \textbf{sem derrubar o site e sem tocar nos dados} (que vivem no banco,
  separados do código).
  \item Mudanças arriscadas podem ser testadas antes em uma \textbf{URL de
  prévia} (cópia isolada que a Vercel gera automaticamente).
  \item Se algo der errado, a Vercel permite \textbf{reverter} para a versão
  anterior em segundos.
\end{enumerate}

\noindent Melhorias comuns são aplicadas à França e replicadas ao Brasil com a
junção das \textit{branches}; cada site se atualiza na sua região.

% ---------------- 11 ----------------
\section{Web Analytics}

O bolão da França usa o \textbf{Web Analytics} da Vercel (gratuito): mostra
número de visitantes, páginas mais acessadas, país e tipo de aparelho. É
\textbf{amigável à privacidade} --- não usa cookies nem identifica indivíduos,
adequado às regras de privacidade da União Europeia. É ativado apenas onde a
variável \texttt{NEXT\_PUBLIC\_ANALYTICS} está definida.

% ---------------- 12 ----------------
\section{Custos}

\begin{center}
\begin{tabularx}{\textwidth}{@{}L >{\bfseries}r@{}}
\toprule
\normalfont Item & Custo \\
\midrule
GitHub (repositório privado) & R\$ 0 \\
Vercel (2 projetos, plano Hobby) & R\$ 0 \\
Neon (2 bancos, plano Free) & R\$ 0 \\
\midrule
\normalfont\bfseries Total mensal & R\$ 0 \\
\bottomrule
\end{tabularx}
\end{center}

\noindent Os planos gratuitos comportam, com folga, dois grupos de 20 a 50
pessoas. O modelo só precisaria de plano pago se o número de usuários crescesse
muito (centenas de milhares de acessos), o que não é o caso.

% ---------------- 13 ----------------
\section{Glossário}

\begin{description}[leftmargin=0pt,style=nextline]
  \item[Frontend] A parte visível: as telas com que o usuário interage.
  \item[Backend] A parte invisível: a lógica de servidor e o acesso ao banco.
  \item[Framework] Conjunto de ferramentas que estrutura o desenvolvimento
  (aqui, o Next.js).
  \item[ORM] Tradutor entre o banco de dados e o código (aqui, o Prisma).
  \item[Server Action] Função de servidor chamada diretamente por um
  formulário, sem uma API separada.
  \item[Deploy] Publicar uma nova versão do aplicativo no ar.
  \item[Branch] Uma ``linha'' de código no Git; usamos uma para cada bolão.
  \item[Variável de ambiente] Configuração definida fora do código (ex.: fuso,
  pontuação), que muda o comportamento sem alterar o programa.
  \item[Serverless] Modelo em que o provedor executa o código sob demanda, sem
  um servidor fixo a administrar.
  \item[JWT] ``Crachá'' digital assinado que comprova a sessão do usuário.
  \item[Hash (bcrypt)] Transformação irreversível da senha, para nunca
  guardá-la em texto puro.
\end{description}

\vfill
\begin{center}
\color{cinza}\small
Club Brésil --- Bolão da Copa 2026 \quad$\bullet$\quad criado por Enzo Vendramin
\end{center}

\end{document}
